
\documentclass[11pt]{article}

% Basic packages
\usepackage{fullpage}
\usepackage{amsmath,amssymb,amsthm}
\usepackage{hyperref}
\usepackage{mathrsfs}
\usepackage{enumitem}
\usepackage{algorithm}
\usepackage{algorithmic}
\usepackage{algpseudocode}
\usepackage{graphicx}
\usepackage{booktabs}
\usepackage{url}

% Theorem environments
\newtheorem{theorem}{Theorem}
\newtheorem{lemma}{Lemma}
\newtheorem{definition}{Definition}

% Title
\title{
{\Large \bf
Succinct and Subquadratic Proofs for Low-Rank Self-Attention \\
via Polynomial Commitments and Freivalds Checking}
}

\author{hyunokoh Oh}
\date{November 2025}

% Authors
\author{
  Hyunok Oh \\
  Hanyang University \\
  \texttt{hoh@hanyang.ac.kr}
}

\date{\today}

\begin{document}
\maketitle

\begin{abstract}
\begin{abstract}
Self-attention is the computational core of modern Transformer models,
but verifying its correctness in zero knowledge remains intractable:
a naïve proof requires encoding an $n\times n$ attention matrix and
performing $O(n^2)$ work inside a SNARK, far beyond practical limits.
In this paper we present the first succinct verification framework for
low-rank self-attention, enabling proof and verification costs that are
independent of the sequence length~$n$.

Our approach leverages the empirical observation that many attention
layers admit a rank-$k$ factorization $A = UV^\top$.  The prover locally
computes the attention output $y = U(V^\top x)$ and proves correctness by
combining (1) polynomial commitments to the columns of $U$ and $V$,  
(2) a Freivalds-style randomized check that reduces matrix--vector
verification to a single scalar identity, and (3) a small SNARK that
binds all committed openings and enforces the final equality.
This commit-and-prove pipeline leaves only an $O(k)$ scalar relation in
the circuit, avoiding any encoding of dense attention.

The resulting protocol achieves prover time $O(nk)$, verifier time
$O(k)$, and proof size $O(k)$, regardless of~$n$.
We further show how the construction naturally extends to multi-head
attention with linear overhead in the number of heads~$h$, while
preserving succinctness.
A prototype implementation confirms the asymptotic behavior and
demonstrates practical feasibility for realistic $n$ and moderate ranks.

Our results introduce the first cryptographically sound and succinct
verification primitive for low-rank attention, opening a new direction
for verifiable Transformer inference.
\end{abstract}

\end{abstract}

\section{Introduction}

\subsection{Motivation}
Self-attention is the computational core of modern Transformer models,
and its quadratic structure poses a fundamental bottleneck for
verifiable machine learning (zkML).  
For a sequence length $n$ and hidden dimension $d$, the attention
mechanism must construct the score matrix $QK^\top \in \mathbb{F}^{n \times n}$,
apply softmax, and project through $V$.
This results in $\Theta(n^2 d)$ arithmetic operations,
and existing verifiable computation frameworks inherit this quadratic cost.

Current zkML systems rely heavily on generic techniques such as
GKR/sum-check protocols~\cite{goldwasser2008gkr} or circuit-based SNARKs,
both of which require arithmetizing the entire $n \times n$ attention matrix.
Even optimized frameworks for verifiable ML (e.g., Hyrax~\cite{zhang2018hyrax})
remain effectively quadratic for self-attention,
making verification impractical for realistic sequence lengths.

A key insight from recent ML research is that the attention matrix
is often \emph{approximately low-rank}.
Empirical studies show that $QK^\top$ can be well-approximated by a
rank-$k$ factorization $U V^\top$ with $k \ll n$, enabling dramatic
reductions in inference cost.
Several models exploit this phenomenon, including
Linformer~\cite{wang2020linformer},
Performer~\cite{choromanski2021performer},
and Nystr\"omformer~\cite{xiong2021nystromformer}.
Despite this, the verification community has not leveraged
low-rank structure for succinct arguments.

If verification could exploit a rank-$k$ factorization,
checking correctness of $y = A x$ for $A = U V^\top$
would require only $O((m+n)k)$ work rather than $\Theta(n^2)$.
Randomized linear-algebra techniques such as Freivalds' algorithm~\cite{freivalds1979fast}
already provide a theoretical blueprint for verifying matrix–vector products
with subquadratic complexity.
However, incorporating these tools into a preprocessing-friendly SNARK pipeline,
while maintaining binding between polynomial commitments and circuit inputs,
requires a nontrivial construction.

This gap motivates our work.
We present a new proof system that combines low-rank structure,
polynomial commitments~\cite{kate2010polynomial},
Freivalds checking, and lightweight SNARK embedding
to eliminate all quadratic structure from the verifier and the circuit.
Our protocol achieves online prover complexity $O((m+n)k)$,
matching the inherent cost of low-rank self-attention,
and provides the first subquadratic verifiable attention mechanism.


\subsection{Challenges in Verifiable Self-Attention}
Designing a succinct proof system for self-attention is challenging due 
to several inherent algebraic and structural obstacles. Existing zkML 
frameworks---including GKR/sum-check based systems~\cite{goldwasser2008gkr} and 
circuit-oriented SNARKs such as Hyrax~\cite{zhang2018hyrax}---struggle to 
handle the quadratic complexity and mixed operations within the 
attention mechanism. We highlight the primary challenges below.

\paragraph{Quadratic Structure of $QK^\top$.}
The core difficulty is that the attention matrix $A = QK^\top$ is 
$n \times n$, requiring $\Theta(n^2 d)$ scalar multiplications.  
Any arithmetization-based proof system must, in principle, process 
each of the $n^2$ bilinear terms unless additional structure is exploited.  
Hence, current zkML systems inherit the quadratic blow-up and cannot 
scale to realistic Transformer sequence lengths.

\paragraph{Heterogeneous Algebraic Operations.}
Self-attention interleaves three operations with fundamentally 
different algebraic properties:
(1) bilinear matrix multiplication $QK^\top$,  
(2) the elementwise softmax nonlinearity, and  
(3) the projection through $V$.  
General-purpose SNARKs represent these with large R1CS or PLONK-style 
constraints, while GKR handles only multilinear arithmetic efficiently.  
Neither approach provides a natural way to jointly optimize these 
operations, leading to deep circuits and poor prover performance.

\paragraph{Commitment Binding Across Domains.}
A verifiable attention computation requires the prover to commit to 
intermediate matrices such as $Q$, $K$, $V$, and possibly $QK^\top$.  
Ensuring that values opened from polynomial commitments (e.g., KZG~\cite{kate2010polynomial}) 
correspond exactly to the SNARK's internal wire assignments requires a 
strong form of input binding.  
Modular SNARK frameworks such as LegoSNARK~\cite{campanelli2019legosnark} support 
this in principle, but integrating commitments across multiple attention 
subcomponents remains nontrivial.

\paragraph{Difficulty Exploiting Low-Rank Structure.}
Although numerous ML works demonstrate that $QK^\top$ often admits a 
rank-$k$ approximation~\cite{wang2020linformer,choromanski2021performer,
xiong2021nystromformer}, no prior verifiable computation framework 
leverages this structure for succinctness.  
GKR and sum-check protocols inherently treat the matrix as dense;  
circuit-based SNARKs require encoding all entries of $U$ and $V$ unless 
specialized gadgets are introduced.  
Thus, existing systems cannot reduce prover complexity below 
$\Theta(n^2)$ despite the underlying low-rank structure.

\paragraph{End-to-End Integration Complexity.}
Even if one could efficiently verify $A x$ for $A = U V^\top$, integrating 
this component into a full Transformer pipeline remains challenging.  
Softmax normalization, multi-head attention, and residual connections 
introduce additional verification costs, and current zkML systems lack 
the modularity to substitute an optimized low-rank proof component 
without redesigning the entire circuit.

These challenges explain why verifiable self-attention has remained an 
open problem, and why no previous system achieves subquadratic prover 
complexity for this fundamental operation.
Our work resolves these obstacles by developing a verification pipeline 
that exploits low-rank structure, relies on Freivalds randomized checking 
for $\Theta((m+n)k)$ soundness amplification~\cite{freivalds1979fast}, and 
uses polynomial commitments with SNARK input binding to efficiently 
compose all components.


\subsection{Our Contributions}

We present the first succinct and subquadratic proof system for
verifying low-rank self-attention.  
Our contributions are summarized as follows.

\begin{itemize}[leftmargin=2em]

  \item \textbf{A new verification framework for low-rank self-attention.}  
  We formally show that when the attention matrix admits a 
  rank-$k$ factorization $A = U V^\top$, the verification of 
  $y = A x$ can be reduced to checking only 
  $O((m+n)k)$ linear operations, rather than the $\Theta(n^2)$ 
  operations required in existing zkML frameworks.
  To our knowledge, this is the first protocol to exploit 
  low-rank structure for succinct verifiable self-attention.

  \item \textbf{A polynomial-commitment–based protocol with 
  preprocessing-friendly structure.}  
  We introduce a construction that commits to $U$ and $V$ using 
  polynomial commitments (e.g., KZG~\cite{kate2010polynomial}) and verifies
  the correctness of all openings while preserving a preprocessing 
  trust model.  
  Our design allows the model owner to commit to $U$ and $V$ once,
  enabling verifiable inference with sublinear verifier effort.

  \item \textbf{A composition of Freivalds checking and lightweight SNARK embedding.}  
  We combine Freivalds' randomized linear-algebra check~\cite{freivalds1979fast}
  with a compact SNARK (e.g., LegoSNARK~\cite{campanelli2019legosnark}) 
  to ensure correct binding between committed values and circuit inputs.
  This eliminates the need to materialize the dense $n \times n$ matrix 
  inside the SNARK, reducing circuit size to $O(k)$.

  \item \textbf{Subquadratic prover time and succinct verifier time.}  
  Our online prover runs in $O((m+n)k)$ time, matching the complexity 
  of low-rank matrix–vector multiplication.
  Verifier time is dominated by a small number of commitment openings 
  and a constant-size SNARK verification, achieving practical succinctness.

  \item \textbf{A verifiable pipeline for Transformer self-attention.}  
  We show how to integrate our low-rank verification component into 
  the full attention pipeline, including the projection through $V$ 
  and the softmax normalization used in Transformer architectures.  
  This establishes the first end-to-end framework for 
  \emph{subquadratic verifiable Transformer inference}.

  \item \textbf{Prototype and empirical evaluation.}  
  We implement a proof-of-concept prototype and provide benchmarks 
  demonstrating that our approach significantly outperforms existing 
  zkML systems such as GKR- or sum-check–based verifiers.  
  In particular, we show concrete reductions in prover time for 
  practical values of $n$ and $k$.
  % (No additional citations needed here.)

\end{itemize}


\subsection{Technical Overview}

We provide a high-level overview of our protocol for verifying 
low-rank self-attention.  
Our goal is to verify, succinctly and soundly, that a prover has 
correctly computed
\[
  y = A x, \qquad A = U V^\top,
\]
where $A \in \mathbb{F}^{m \times n}$ is a rank-$k$ operator derived 
from a self-attention layer.  
The verifier holds commitments to $U$ and $V$ and wishes to check 
the correctness of $y$ without processing the dense matrix $A$.

\vspace{0.7em}
\paragraph{Preprocessing: committing to $U$ and $V$.}
In the preprocessing phase, the model owner commits to the matrices 
$U \in \mathbb{F}^{m \times k}$ and $V \in \mathbb{F}^{n \times k}$ 
using polynomial commitments such as KZG~\cite{kate2010polynomial}.
Each column of $U$ and $V$ is encoded as a polynomial, 
and the verifier stores the resulting commitments.
This allows the prover to later open evaluations of $U$ and $V$ at 
verifier-chosen points while preventing equivocation.
This preprocessing reflects realistic inference settings in which the 
attention parameters remain fixed across many queries.

\vspace{0.7em}
\paragraph{Online computation: factorized matvec.}
Given an input vector $x \in \mathbb{F}^n$, the prover computes
\[
  z = V^\top x \qquad (\text{cost } O(nk)),
\]
and then
\[
  y = U z \qquad (\text{cost } O(mk)).
\]
This reproduces the effect of multiplying by the 
rank-$k$ matrix $A = U V^\top$ while avoiding the 
$\Theta(n^2)$ cost of explicitly forming $A$.

\vspace{0.7em}
\paragraph{Randomized verification via Freivalds' check.}
To verify the relation $y = A x$, the verifier samples a random 
challenge vector $r \leftarrow \mathbb{F}^m$ and checks whether
\[
  r^\top y = (r^\top U)(V^\top x).
\]
This reduces verification to evaluating:
\[
  s_1 = r^\top y, \qquad 
  s_2 = u^\top z \quad \text{where } u = r^\top U.
\]
By Freivalds' classical result~\cite{freivalds1979fast},
the equality $s_1 = s_2$ holds with probability at least 
$1 - 1/|\mathbb{F}|$ if and only if $y = A x$.
Repeating with a few random challenges yields overwhelming soundness 
with negligible verifier cost.

\vspace{0.7em}
\paragraph{Polynomial openings and input binding.}
To compute $u = r^\top U$, the prover must open the committed columns 
of $U$ at the verifier-chosen evaluation points.
Similarly, to confirm correctness of $z = V^\top x$, 
the prover opens the necessary components of $V$.
These openings are accompanied by a KZG proof~\cite{kate2010polynomial} or similar, 
allowing the verifier to check correctness in logarithmic time.

A key challenge is ensuring that the openings provided to the verifier 
match the values supplied inside the SNARK circuit.  
We use a lightweight SNARK such as LegoSNARK~\cite{campanelli2019legosnark} to 
\emph{bind} the polynomial openings to the arithmetic circuit that checks 
the equality $s_1 = s_2$.
This guarantees that the prover cannot equivocate by providing 
inconsistent values across different components of the protocol.

\vspace{0.7em}
\paragraph{Circuit size: $O(k)$, independent of $n$.}
The SNARK circuit merely checks the scalar equality 
$r^\top y = u^\top z$, which requires $O(k)$ multiplications and 
a few additions.
No part of the $n \times n$ matrix $A$ is ever materialized or encoded 
in the circuit.
Thus the SNARK is \emph{constant-sized with respect to the sequence 
length~$n$}, providing genuine succinctness.

\vspace{0.7em}
\paragraph{Extension to self-attention.}
In the case of Transformers, the matrices $U$ and $V$ arise from 
a low-rank approximation of the attention scores $QK^\top$, as observed 
in Linformer~\cite{wang2020linformer}, Performer~\cite{choromanski2021performer}, 
and related architectures.
Our protocol verifies the projection step 
$y = \mathrm{softmax}(QK^\top) V$ by:
(1) verifying $U V^\top x$ as above, and  
(2) incorporating a lightweight gadget for enforcing the softmax 
normalization and attention weights.
Since these steps do not require encoding the full attention matrix, 
our approach yields the first end-to-end \emph{subquadratic} 
verifiable self-attention pipeline.



\section{Background}

\subsection{Self-Attention and Low-Rank Approximations}

Self-attention computes interactions among all token embeddings in a 
sequence and is central to Transformer architectures.  
Given input vectors 
$Q, K, V \in \mathbb{F}^{n \times d}$,  
the attention mechanism constructs the score matrix
\[
  S = \frac{Q K^\top}{\sqrt{d}},
\]
applies row-wise softmax to obtain attention weights,
\[
  A = \mathrm{softmax}(S),
\]
and computes the output
\[
  Y = A V.
\]
The computational bottleneck arises from forming the $n \times n$ matrix 
$S = Q K^\top$, which requires $\Theta(n^2 d)$ operations, and from 
subsequent multiplication with $V$.
Any verifiable computation framework that attempts to arithmetize 
$A$ or $S$ as dense matrices inevitably inherits this quadratic cost.

\vspace{0.7em}
\paragraph{Empirical low-rank structure of attention.}
Recent empirical analyses reveal that attention matrices 
often exhibit strong low-rank structure.  
Several architectures exploit this finding to reduce inference cost:  
Linformer projects $K$ and $V$ into a lower-dimensional subspace, 
arguing that attention scores lie approximately in a 
rank-$k$ subspace with $k \ll n$~\cite{wang2020linformer}.  
Performer models attention via random feature maps, showing that 
$Q K^\top$ can be approximated by a kernelized low-rank factorization 
using FAVOR+ features~\cite{choromanski2021performer}.  
Nystr\"omformer constructs a Nystr\"om approximation of attention, 
explicitly modeling $A$ as a low-rank matrix derived from a small set of 
landmark points~\cite{xiong2021nystromformer}.  
These works collectively demonstrate that the effective rank of 
self-attention is often far smaller than $n$, especially in long-context settings.

\vspace{0.7em}
\paragraph{Low-rank modeling via factorization.}
A rank-$k$ approximation of $QK^\top$ can be written as
\[
  Q K^\top \approx U V^\top,
\]
where $U, V \in \mathbb{F}^{n \times k}$.
Such factorizations arise naturally from:
(i) projection-based methods (e.g., Linformer),  
(ii) kernelized linearization techniques (e.g., Performer), and  
(iii) Nystr\"om or landmark-based matrix approximations.  
In all cases, the approximation captures the dominant subspace of 
attention scores while reducing computational complexity to $O(n k)$.

\vspace{0.7em}
\paragraph{Benefits for verifiable computation.}
Crucially, a rank-$k$ factorization transforms the quadratic 
matrix--vector multiplication
\[
  y = (QK^\top) x
\]
into
\[
  y = U (V^\top x),
\]
which requires $O(n k)$ operations.  
This matches the cost of low-rank inference in ML systems, and it 
enables subquadratic verification when combined with 
randomized linear-algebra techniques (e.g., Freivalds checking) 
and succinct proofs.

Thus, low-rank approximation provides both a 
\emph{modeling advantage} used in efficient Transformers and a 
\emph{verification advantage} leveraged in our proof system.


\subsection{Polynomial Commitments (KZG/IPA)}

Polynomial commitment schemes enable a prover to commit to a polynomial
$f(X)$ and later open evaluations $f(\alpha)$ at verifier-chosen points
with succinct proofs.  
They are a fundamental building block in modern zk-SNARK constructions
and form the backbone of our mechanism for committing to the low-rank
matrices $U$ and $V$.

\paragraph{KZG commitments.}
The pairing-based KZG scheme of Kate, Zaverucha, and Goldberg~\cite{kate2010polynomial}
commits to a polynomial by embedding its evaluation at a secret point $s$
into a group generator:
\[
  \mathsf{cm} = g^{f(s)}.
\]
To open $f(\alpha)$, the prover constructs a quotient polynomial
\[
  w(X) = \frac{f(X) - f(\alpha)}{X - \alpha}
\]
and sends the proof $\pi = g^{w(s)}$.
The verifier checks the pairing equation
\[
  e(\mathsf{cm} / g^{f(\alpha)}, g)
  \stackrel{?}{=} 
  e(\pi, g^{s - \alpha}),
\]
providing constant-size proofs and $O(1)$ verification.
KZG commitments are ideal for preprocessing-heavy settings such as
verifiable inference, where $U$ and $V$ remain fixed across queries.

\paragraph{Transparent IPA-based commitments.}
In environments without pairings or trusted setup, one may use
inner-product–argument (IPA)–based polynomial commitments.
The DARK framework~\cite{bunz2019dark} gave the first fully transparent
PCS using recursive inner-product arguments, achieving logarithmic-size
proofs and no trusted setup.
This line of work was further improved by the efficient transparent
scheme of Kim et al.~\cite{kim2023etpc}, which reduces prover time,
proof size, and verification overhead while maintaining a fully
transparent setup.
These schemes provide attractive alternatives to KZG in pairing-free or
public-parameter–sensitive settings.

\paragraph{Why polynomial commitments are needed here.}
In our construction, each column of $U$ and $V$ is encoded as a polynomial.
This allows the verifier to request openings of specific entries
(e.g., those needed to compute $u = r^\top U$ during Freivalds' check)
without requiring the prover to reveal entire columns.
Succinct opening proofs ensure that the prover cannot equivocate:
the same matrix entry must be used consistently across the online
computation, the Freivalds relation, and the SNARK circuit.

\paragraph{SNARK input binding.}
To prevent mismatches between polynomial openings and circuit inputs,  
we embed these scalar openings inside a lightweight SNARK such as  
LegoSNARK~\cite{campanelli2019legosnark}.
This ensures that the prover cannot supply one value to the commitment
verifier and a different value to the SNARK circuit checking
\[
  r^\top y = u^\top z.
\]
Thus, polynomial commitments act as the cryptographic glue connecting
the off-chain low-rank algebra to the succinct arithmetic circuit.

\paragraph{Summary.}
Polynomial commitments---whether pairing-based (KZG) or transparent
IPA-based (DARK, Kim et al. 2023)---are essential for enabling
our preprocessing-friendly, subquadratic proof system.
They allow the verifier to check only a handful of openings while
delegating all heavy algebra to the prover.

\subsection{Freivalds Randomized Checking}

Freivalds' algorithm~\cite{freivalds1979fast} provides a simple and 
powerful randomized method for verifying matrix–vector and 
matrix–matrix multiplications.
It guarantees correctness with high probability while requiring 
time only linear in the size of the operands, making it ideally 
suited for verifying low-rank computations.

\paragraph{Classical Freivalds check.}
Given matrices $A, B$ and a claimed product $C$, Freivalds' check 
verifies $AB = C$ by sampling a random vector 
$r \leftarrow \mathbb{F}^n$ and checking whether
\[
  A (B r) = C r.
\]
If $AB = C$, the equality always holds.
If $AB \neq C$, then the probability (over $r$) that the equality 
nevertheless holds is at most $1/|\mathbb{F}|$, regardless of the size 
of the matrices.
Thus, with $t$ repetitions, the soundness error becomes $\leq 
(|\mathbb{F}|^{-1})^t$, negligible in cryptographic fields.

\paragraph{Adaptation to matrix–vector products.}
In our setting, we want to verify the correctness of 
\[
  y = A x
\]
for the rank-$k$ matrix $A = U V^\top$.
Instead of computing the full matrix, the prover computes
\[
  z = V^\top x, \qquad
  y = U z.
\]
The verifier samples a random $r \in \mathbb{F}^m$ and checks
\[
  r^\top y \stackrel{?}{=} (r^\top U)(V^\top x).
\]
Define $u = r^\top U$.
The right-hand side becomes
\[
  u^\top z,
\]
so the check reduces to comparing two scalars 
$s_1 = r^\top y$ and $s_2 = u^\top z$.

If $y = U(V^\top x)$, the equality holds identically.
If not, Freivalds' bound shows
\[
  \Pr[s_1 = s_2] \leq \frac{1}{|\mathbb{F}|}.
\]
Hence, a single check already provides strong probabilistic soundness.

\paragraph{Why Freivalds works especially well for low-rank structure.}
A key advantage of the low-rank factorization is that 
Freivalds' algorithm interacts \emph{only} with the compressed 
representations $U$ and $V^\top x$, not with the full $n \times n$ 
matrix $A$.
All quantities needed for verification are linear forms:
\[
  r^\top y,\quad r^\top U,\quad V^\top x.
\]
These can be committed and opened succinctly using polynomial 
commitments, and checked inside a very small SNARK circuit.
Thus Freivalds' check allows verification complexity to scale as 
$O((m+n)k)$, matching the true algebraic complexity of the low-rank 
matvec itself.

\paragraph{Embedding Freivalds in SNARKs.}
Freivalds' algorithm produces a single scalar equation
\[
  r^\top y = (r^\top U)(V^\top x),
\]
which is inexpensive to encode inside a SNARK.
The circuit size is $O(k)$ because:
(i) $u = r^\top U$ consists of $k$ multiplications,  
(ii) $u^\top z$ adds another $k$ multiplications, and  
(iii) the final comparison is constant-sized.

Crucially, the commitment openings for $U$ and $V$ are tied to these 
circuit inputs via an input-binding SNARK (e.g., LegoSNARK), ensuring 
that the prover cannot cheat by providing inconsistent values.

\paragraph{Summary.}
Freivalds' randomized linear-algebra check provides the key tool that 
removes the quadratic bottleneck of self-attention verification.
It turns the $n \times n$ structure of $A = U V^\top$ into a single 
scalar equation whose correctness can be enforced succinctly.
Combined with polynomial commitments and SNARK binding, this yields a 
subquadratic and preprocessing-friendly verification pipeline.

\subsection{SNARKs and Input Binding}

Polynomial commitments alone guarantee that individual openings are
correct with respect to a committed polynomial, but they do not ensure
that the \emph{same values} are used consistently across different
components of the protocol.
In our setting, the prover must supply:
(i) openings of selected entries of $U$ and $V$,  
(ii) the vectors $z = V^\top x$ and $y = U z$, and  
(iii) the scalar values used in the Freivalds equation
\[
  r^\top y = (r^\top U)(V^\top x).
\]
Without a mechanism for enforcing consistency among these values, a
dishonest prover could satisfy the commitment verifier and the arithmetic
circuit separately by providing different values to each component.

\paragraph{The need for input binding.}
The verifier checks polynomial openings using KZG or IPA-based PCS,
while the SNARK circuit checks the scalar relation arising from
Freivalds' check.
To prevent the prover from equivocating between these two layers,
we require a binding mechanism that cryptographically links the
openings to the circuit inputs.
This concept, known as \emph{input binding}, ensures that the data
originating from polynomial commitments is faithfully reflected inside
the SNARK.

\paragraph{LegoSNARK as a modular input-binding framework.}
Modern composable SNARK frameworks such as LegoSNARK~\cite{campanelli2019legosnark}
explicitly support modular composition of external commitments with
internal circuit logic.
LegoSNARK provides a gadget that allows the circuit to assert that a
publicly supplied value is consistent with an opening verified outside
the circuit.
Concretely, the SNARK enforces an equality of the form
\[
  v_\text{circuit} = v_\text{commit},
\]
where $v_\text{commit}$ is the field element obtained from the
polynomial commitment opening.
This prevents the prover from supplying inconsistent values in the two
parts of the protocol.

\paragraph{Binding the Freivalds relation.}
For the Freivalds equation
\[
  r^\top y = u^\top z,
\]
the SNARK circuit requires the entries of $u = r^\top U$ and $z =
V^\top x$ to be correct.
These entries are obtained from polynomial openings of $U$ and $V$,
respectively.
Through the input-binding gadget, each such entry is imported into the
circuit in a way that prevents substitution or modification.
Thus the SNARK enforces that the same values used to verify polynomial
openings are precisely the values used to check the scalar equation.

\paragraph{Minimal circuit overhead.}
A key advantage of input binding in this context is that the required
SNARK circuitry remains extremely small.
Once the circuit receives the bound inputs, it only implements:
(i) $k$ multiplications to compute $u^\top z$,  
(ii) $k$ multiplications for $r^\top y$, and  
(iii) a final equality check.
No reconstruction of the matrix $A = U V^\top$ or of the vectors
$U z$ or $V^\top x$ is needed inside the circuit.
Thus the binding mechanism ensures soundness without introducing
quadratic overhead.

\paragraph{Summary.}
Input binding is indispensable for ensuring the soundness of our
protocol.
It ties polynomial commitment openings to the SNARK circuit, preventing
equivocation and enabling a modular design where each component can be
verified independently but still interact securely.
This binding is what makes it possible to combine Freivalds' randomized
check, polynomial commitments, and succinct proving into a unified,
subquadratic verification framework.


\section{Problem Formulation}

We formalize the setting of verifiable low-rank self-attention.
Our goal is to enable a computationally powerful prover to convince a
resource-limited verifier that a self-attention layer was evaluated
correctly, without requiring the verifier to handle any $n \times n$
intermediate matrices.

\subsection{Computation Model}

Let $n$ denote the sequence length and $d$ the embedding dimension.
A self-attention layer computes:
\[
  S = \frac{QK^\top}{\sqrt{d}}, \qquad
  A = \mathrm{softmax}(S), \qquad
  Y = A V,
\]
where $Q, K, V \in \mathbb{F}^{n \times d}$.
Direct verification of $A$ or $S$ requires operating on dense
$n \times n$ matrices, which is infeasible for large $n$.

Following prior work on efficient attention approximations, we assume the
existence of a rank-$k$ factorization
\[
  A \approx U V^\top,
\]
where $U, V \in \mathbb{F}^{n \times k}$ with $k \ll n$.
The target computation becomes
\[
  y = A x \approx U(V^\top x),
\]
for some input vector $x \in \mathbb{F}^n$.

\subsection{Preprocessing Assumptions}

We consider a preprocessing model where the matrices $U$ and $V$ are
public parameters that remain fixed across many proofs.
During preprocessing, the model owner commits to $U$ and $V$ using a
polynomial commitment scheme (e.g., KZG or IPA-based PCS).
The commitments are public and the prover cannot later alter $U$ or $V$.

This reflects realistic inference deployments: once a Transformer layer
is fixed, it is reused across many queries.

\subsection{Prover and Verifier Capabilities}

\paragraph{Prover.}
The prover:
\begin{itemize}
  \item receives an input vector $x$,
  \item computes $z = V^\top x$ and $y = U z$,
  \item proves that $y$ was computed correctly with respect to the
        committed matrices $U, V$.
\end{itemize}
The prover is computationally unbounded except for polynomial-time
constraints.

\paragraph{Verifier.}
The verifier:
\begin{itemize}
  \item holds commitments to $U$ and $V$,
  \item can request polynomial openings at selected indices,
  \item generates random challenges $r$ for Freivalds' check,
  \item verifies a succinct SNARK proof for the equation
        \[
          r^\top y = (r^\top U)(V^\top x),
        \]
  \item must run in time sublinear in $n$ (ideally $O(k)$ or constant).
\end{itemize}

\subsection{Threat Model}

We assume a malicious prover who may try to:
\begin{itemize}
  \item produce incorrect outputs $y$ inconsistent with $U, V$,
  \item forge polynomial openings for $U$ or $V$,
  \item supply inconsistent values to the commitment verifier and to the
        SNARK circuit (equivocation),
  \item exploit low-rank structure to fabricate “shortcut’’ outputs not
        derived from the intended computation.
\end{itemize}

The verifier must detect all such deviations except with negligible
probability.

\subsection{Security Goals}

Our proof system must satisfy the following properties:

\paragraph{Completeness.}
If the prover computes $y = U(V^\top x)$ honestly, the verifier must
accept with probability~1.

\paragraph{Soundness.}
If $y \neq U(V^\top x)$, then for a random challenge
$r \leftarrow \mathbb{F}^n$, the probability that the prover convinces
the verifier is at most $1/|\mathbb{F}|$ (for a single Freivalds check),
and negligible after a small number of repetitions.

\paragraph{Input binding.}
Values used inside the SNARK must match the values obtained from
polynomial openings of $U$ and $V$.
The prover must not be able to equivocate between these sources.

\paragraph{Succinctness.}
The verifier must run in time sublinear in $n$ and the SNARK proof size
must be independent of $n$ (ideally $O(k)$ or constant).

\paragraph{Preprocessing soundness.}
Once commitments to $U$ and $V$ are fixed during preprocessing, the
prover cannot alter them or switch to alternative matrices in later
proofs.

\subsection{Problem Definition}

Given:
\begin{itemize}
  \item commitments to $U, V \in \mathbb{F}^{n \times k}$,
  \item an input vector $x \in \mathbb{F}^n$,
  \item an output vector $y \in \mathbb{F}^n$,
\end{itemize}
the prover must convince the verifier that
\[
  y = U (V^\top x),
\]
or equivalently,
\[
  r^\top y = (r^\top U)(V^\top x)
\]
for a verifier-chosen random $r$,
with negligible soundness error and with verifier time independent of
$n$.

This captures the precise correctness requirement for verifiable
low-rank self-attention.

\section{Protocol Overview}

We give a high-level description of our proof system for verifying 
low-rank self-attention.  
The protocol consists of four interacting components:
(1) preprocessing commitments to the low-rank matrices $U$ and $V$,
(2) online computation of the low-rank matvec $y = U(V^\top x)$,
(3) a Freivalds-based randomized linear check, and
(4) a succinct SNARK that binds all components together.

\subsection{1. Preprocessing: Committing to $U$ and $V$}

Before any proofs are generated, the model owner commits to the matrices
$U, V \in \mathbb{F}^{n \times k}$ using a polynomial commitment 
scheme (KZG or IPA-based).
Each column is encoded as a polynomial and committed once.
These commitments act as immutable public parameters for all future
queries.

This preprocessing captures the practical setting of inference:
the self-attention layer is fixed, while inputs $x$ vary across tasks.

\subsection{2. Online Computation: Computing $z$ and $y$}

Given an input vector $x \in \mathbb{F}^n$, the prover computes:
\[
  z = V^\top x \qquad (\text{cost } O(nk)),
\]
\[
  y = U z \qquad (\text{cost } O(nk)).
\]
No dense $n \times n$ matrix is ever formed.
The prover produces $y$ as the claimed output of the self-attention
projection.

\subsection{3. Randomized Verification via Freivalds' Check}

To verify correctness of $y$, the verifier samples a random challenge
$r \leftarrow \mathbb{F}^n$ and checks:
\[
  r^\top y \stackrel{?}{=} (r^\top U)(V^\top x).
\]
This reduces verification to a single scalar equality.
To compute the right-hand side, the prover must open selected entries of
$U$ and $V$ at positions determined by $r$ and $x$.
These openings are accompanied by polynomial commitment proofs
attesting to their correctness.

Freivalds' guarantee ensures that if $y \neq U(V^\top x)$, the probability
that the equality holds is at most $1/|\mathbb{F}|$, negligible after
only a few repetitions.

\subsection{4. SNARK for Input Binding and Scalar Verification}

Finally, a lightweight SNARK circuit enforces that all scalar values
supplied to the Freivalds check are consistent with the polynomial
openings.
The circuit takes as inputs:
\[
  r^\top y,\quad u = r^\top U,\quad z = V^\top x,
\]
and verifies the scalar equation:
\[
  r^\top y = u^\top z.
\]

An input-binding gadget (e.g., via LegoSNARK) ensures that these values
match the entries opened from commitments to $U$ and $V$.
Since the circuit encodes only $O(k)$ multiplications, the SNARK proof
is independent of $n$.

\subsection{Putting the Components Together}

Overall, the protocol proceeds as follows:
\begin{enumerate}
  \item \textbf{Setup:} Commit to $U$ and $V$ using polynomial commitments.
  \item \textbf{Prove:}  
    \begin{enumerate}
      \item Compute $z = V^\top x$, $y = U z$.
      \item Respond to polynomial-opening queries for entries of $U$ and $V$.
      \item Construct a SNARK proof for the Freivalds scalar equation.
    \end{enumerate}
  \item \textbf{Verify:}
    \begin{enumerate}
      \item Sample random $r$ and request the necessary openings.
      \item Verify the polynomial openings.
      \item Verify the SNARK proof for the scalar relation.
    \end{enumerate}
\end{enumerate}

The combination of preprocessing, low-rank computation, randomized
linear algebra, and succinct input-binding yields a verification process
whose cost is independent of $n$ and depends only on the rank parameter
$k$.

\section{Construction}

We now describe our full proof system for verifiable low-rank
self-attention.  We begin by specifying the preprocessing setup that
produces the public parameters and commitments used across all
queries.

\subsection{Setup}

The setup phase is executed once by the model owner and produces
public parameters that will be reused for all subsequent proofs.
It consists of two parts:
(i) generating the polynomial-commitment parameters, and
(ii) committing to the fixed low-rank matrices $U$ and $V$.

\paragraph{Polynomial-commitment parameters.}
Depending on the chosen instantiation, the setup proceeds as follows:

\begin{itemize}
  \item \textbf{KZG-based setup.}
  A structured reference string (SRS) is generated for KZG
  commitments~\cite{kate2010polynomial}.
  The SRS includes group elements
  $\{g^{s^i}\}_{i=0}^{\ell}$ for a secret $s$ and degree bound $\ell$.
  This setup is trusted but produces constant-size openings 
  and $O(1)$ verification.

  \item \textbf{Transparent IPA-based setup.}
  A transparent parameter set for an IPA polynomial commitment scheme
  (e.g., DARK~\cite{bunz2019dark} or the improved construction of
  Kim et al.~\cite{kim2023etpc}) is generated.
  No trusted trapdoor is required, at the cost of logarithmic proof size.
\end{itemize}

Let $\mathsf{pp}_{\mathsf{PC}}$ denote the public parameters for the 
chosen polynomial-commitment scheme.

\paragraph{Matrix encoding and commitment.}
Let $U, V \in \mathbb{F}^{n \times k}$ be the low-rank matrices used to
approximate the self-attention operator.
For each column $U_j$ of $U$, the model owner defines a polynomial
$f_{U_j}(X)$ of degree $< n$ such that
\[
  f_{U_j}(i) = U_{i,j} \qquad \text{for all } i \in \{1,\dots,n\}.
\]
Similarly, each column $V_j$ of $V$ is encoded as a polynomial
$f_{V_j}(X)$.

Using $\mathsf{pp}_{\mathsf{PC}}$, the model owner computes the
commitments:
\[
  \mathsf{cm}_{U_j} = \mathsf{Commit}(f_{U_j}), \qquad
  \mathsf{cm}_{V_j} = \mathsf{Commit}(f_{V_j}),
\]
for every $j \in \{1,\dots,k\}$.

Let
\[
  \mathsf{CM}_U = (\mathsf{cm}_{U_1}, \dots, \mathsf{cm}_{U_k}), \qquad
  \mathsf{CM}_V = (\mathsf{cm}_{V_1}, \dots, \mathsf{cm}_{V_k})
\]
denote the full set of commitments.

\paragraph{Published public parameters.}
The following values become public and fixed for all proof sessions:
\begin{itemize}
  \item the polynomial-commitment parameters $\mathsf{pp}_{\mathsf{PC}}$,
  \item the commitments $\mathsf{CM}_U$ and $\mathsf{CM}_V$,
  \item the rank parameter $k$ and the length parameter $n$.
\end{itemize}

\paragraph{Security requirement.}
After setup, the prover must not be able to modify $U$ or $V$.
All openings in later phases must correspond to the polynomials committed
during this setup.

\paragraph{Output of Setup.}
The setup algorithm outputs the public parameters:
\[
  \mathsf{pp} = (\mathsf{pp}_{\mathsf{PC}},\, \mathsf{CM}_U,\, 
                 \mathsf{CM}_V,\, n,\, k).
\]
These parameters are given to the verifier and to all provers.
No further trusted actions are required.
\subsection{Prover Algorithm}

Given the public parameters 
$\mathsf{pp} = (\mathsf{pp}_{\mathsf{PC}}, \mathsf{CM}_U, \mathsf{CM}_V, n, k)$
and an input vector $x \in \mathbb{F}^n$, the prover must compute the
claimed output $y$ and produce a proof that certifies the correctness of
the relation
\[
  y = U(V^\top x)
\]
with respect to the committed matrices $U$ and $V$.
The prover algorithm consists of three stages:
(1) online computation,  
(2) responding to polynomial-opening queries,  
and (3) constructing a SNARK proof for Freivalds’ relation.

\subsubsection{1. Online Computation}

The prover computes the low-rank matvec using the factorization
$A = U V^\top$:
\[
  z = V^\top x  \qquad\text{and}\qquad  y = U z.
\]
These operations require only $O(nk)$ time and do not involve materializing
the full $n \times n$ matrix $A$.

The prover outputs $y$ as the claimed result.

\subsubsection{2. Opening Queries for $U$ and $V$}

After receiving a random challenge $r \leftarrow \mathbb{F}^n$ from the
verifier (as part of Freivalds' check), the prover must provide openings
for the entries needed to compute:
\[
  u = r^\top U, \qquad\quad z = V^\top x.
\]

\paragraph{Openings from $U$.}
To compute $u = r^\top U$, the prover must reveal the entries
\[
  \{ U_{i,j} : i \in \mathrm{supp}(r),\ j \in [k] \},
\]
along with polynomial-opening proofs that these entries match the 
committed polynomials $f_{U_j}$.

\paragraph{Openings from $V$.}
Similarly, to justify the computation of $z = V^\top x$, the prover
reveals:
\[
  \{ V_{i,j} : i \in \mathrm{supp}(x),\ j \in [k] \},
\]
again with matching polynomial-opening proofs.

Let all revealed scalar values be collected into a vector:
\[
  \mathsf{openings} = \{u, z, \text{openings from } U, \text{openings from } V\}.
\]

\subsubsection{3. Constructing the SNARK Proof}

The prover now builds a SNARK proof $\pi$ that certifies the scalar
relation
\[
  r^\top y = u^\top z.
\]

The SNARK circuit enforces:
\begin{enumerate}
  \item the equality $s_1 = r^\top y$ (public input),
  \item the equality $s_2 = u^\top z$ (public input),
  \item the final relation $s_1 = s_2$,
  \item input-binding constraints ensuring that all values 
        $(u, z)$ correspond exactly to the polynomial openings.
\end{enumerate}

A modular SNARK such as LegoSNARK provides the necessary input-binding
gadget to tie $\mathsf{openings}$ to the circuit inputs.
The circuit size is $O(k)$, as it includes only $k$ multiplications to
compute each inner product.

The prover outputs:
\[
  \pi_{\mathrm{PC}} \;=\; \text{all polynomial-opening proofs},\qquad
  \pi_{\mathrm{SNARK}} \;=\; \text{the SNARK proof}.
\]

\subsubsection{Prover Summary}

The prover outputs the tuple:
\[
  (\; y,\, \pi_{\mathrm{PC}},\, \pi_{\mathrm{SNARK}} \;).
\]

\begin{algorithm}[t]
\caption{Prover Algorithm}
\begin{algorithmic}[1]
\Require{$x$, public parameters $\mathsf{pp}$, commitments $\mathsf{CM}_U,\mathsf{CM}_V$}
\State Compute $z \gets V^\top x$
\State Compute $y \gets U z$
\State Receive random challenge $r$ from verifier
\State Compute $u \gets r^\top U$
\State Generate polynomial openings for required entries of $U$ and $V$
\State Build SNARK witness for $r^\top y = u^\top z$
\State Construct SNARK proof $\pi_{\mathrm{SNARK}}$
\State \Return $(y,\; \pi_{\mathrm{PC}},\; \pi_{\mathrm{SNARK}})$
\end{algorithmic}
\end{algorithm}

\subsection{Verifier Algorithm}

Given public parameters 
$\mathsf{pp} = (\mathsf{pp}_{\mathsf{PC}}, \mathsf{CM}_U, \mathsf{CM}_V, n, k)$ 
and a prover-generated proof 
$(y,\, \pi_{\mathrm{PC}},\, \pi_{\mathrm{SNARK}})$,
the verifier must check that the claimed output $y$ satisfies
\[
  y = U(V^\top x),
\]
with respect to the committed matrices $U$ and $V$.
The verifier proceeds in three steps:
(1) randomized challenge generation,  
(2) verification of polynomial-opening proofs, and  
(3) verification of the SNARK proof.

\subsubsection{1. Sampling a Random Challenge}

The verifier samples a random vector
\[
  r \leftarrow \mathbb{F}^n
\]
to instantiate Freivalds' check.
The vector $r$ is sent to the prover, who must open the corresponding
entries of $U$ and $V$.

To avoid handling full-length vectors, the verifier may sample $r$ from
any distribution with full support (e.g., sparse random vectors), as
Freivalds' correctness requires only that $r$ has at least one random
coordinate.

\subsubsection{2. Verifying Polynomial-Opening Proofs}

The prover returns a set of polynomial openings and opening proofs 
$\pi_{\mathrm{PC}}$ for the entries used to compute:
\[
  u = r^\top U, 
  \qquad
  z = V^\top x.
\]

For each opened value $v_{i,j}$, the verifier checks:
\[
  \mathsf{VerifyOpen}(\mathsf{cm}_{U_j},\, i,\, v_{i,j}) = 1,
\]
or the analogous check for entries of $V$.

If any polynomial-opening proof fails, the verifier rejects.

\subsubsection{3. Verifying the SNARK Proof}

To confirm correctness of the scalar equation
\[
  r^\top y = (r^\top U)(V^\top x),
\]
the verifier checks the SNARK proof $\pi_{\mathrm{SNARK}}$ containing:

\begin{itemize}
  \item The public inputs: $r$, $y$, and all opened values needed to 
        compute $u$ and $z$.
  \item The scalar values $s_1 = r^\top y$ and $s_2 = u^\top z$.
  \item A constraint enforcing $s_1 = s_2$.
  \item Input-binding constraints ensuring that $u$ and $z$ match the 
        corresponding polynomial openings verified in the previous step.
\end{itemize}

The verifier runs the SNARK verifier for $\pi_{\mathrm{SNARK}}$, which 
takes time $O(1)$ or $O(k)$ depending on the instantiation.

\subsubsection{Verifier Decision}

The verifier accepts if and only if:
\begin{enumerate}
  \item all polynomial-opening proofs in $\pi_{\mathrm{PC}}$ verify, and
  \item the SNARK proof $\pi_{\mathrm{SNARK}}$ is valid.
\end{enumerate}

Otherwise, the verifier outputs \emph{reject}.

\begin{algorithm}[t]
\caption{Verifier Algorithm}
\begin{algorithmic}[1]
\Require{public parameters $\mathsf{pp}$, input $x$, proof $(y,\pi_{\mathrm{PC}},\pi_{\mathrm{SNARK}})$}
\State Sample random $r \leftarrow \mathbb{F}^n$
\State Send $r$ to prover
\State Receive polynomial openings and proofs $\pi_{\mathrm{PC}}$
\If{any $\mathsf{VerifyOpen}$ fails} \Return reject
\EndIf
\If{$\mathsf{VerifySNARK}(\pi_{\mathrm{SNARK}}) = 1$}
    \Return accept
\Else
    \Return reject
\EndIf
\end{algorithmic}
\end{algorithm}

\paragraph{Verifier Cost.}
The verifier’s cost consists of:
\[
  O(\text{polylog}(n)) \text{ for polynomial openings} \quad+\quad O(k)
  \text{ for SNARK verification}.
\]
Thus the verification time is independent of $n$ and depends only on the
rank parameter $k$.

\subsection{Proof Structure}

We describe the structure of the proof generated by the prover and the
values that must be provided to the verifier.  
A complete proof consists of:
(1) the claimed output $y$,  
(2) polynomial-opening proofs for all queried entries of $U$ and $V$,  
and (3) a succinct SNARK proof that certifies the scalar Freivalds relation.

\subsubsection{Public Inputs}

The verifier receives the following as public inputs:

\begin{itemize}
  \item the input vector $x$,
  \item the claimed output vector $y$,
  \item the random challenge $r$ sampled by the verifier,
  \item the polynomial commitments $\mathsf{CM}_U, \mathsf{CM}_V$,
  \item the public parameters $\mathsf{pp}_{\mathsf{PC}}$,
  \item the scalar values
        \[
          s_1 = r^\top y,\qquad s_2 = u^\top z,
        \]
        where $u = r^\top U$ and $z = V^\top x$,
  \item all field elements revealed during polynomial openings.
\end{itemize}

These values allow the verifier to evaluate the SNARK circuit and
validate the correctness of the polynomial-opening proofs.

\subsubsection{SNARK Witness}

The prover’s private witness consists of:

\begin{itemize}
  \item the vectors $u = r^\top U$ and $z = V^\top x$,
  \item any intermediate field elements used to compute the inner products,
  \item the binding relations that map each scalar in the circuit
        to the corresponding polynomial-opening value,
  \item auxiliary data required by the polynomial-commitment scheme
        (quotient polynomials, witnesses, etc., but not revealed externally).
\end{itemize}

The witness is never revealed; it is used only to satisfy the circuit.

\subsubsection{Circuit Constraints}

The SNARK circuit contains the following constraints:

\paragraph{Inner-product constraints.}
The circuit computes:

\[
  s_1 = r^\top y = \sum_{i=1}^n r_i y_i,
\]
\[
  s_2 = u^\top z = \sum_{j=1}^k u_j z_j.
\]

As $k \ll n$, only the computation of $s_2$ incurs circuit overhead
($O(k)$ multiplications).
The value $s_1$ is given as a public input and does not require
recomputation inside the circuit.

\paragraph{Final correctness constraint.}

\[
  s_1 = s_2.
\]

This single scalar equation represents the full Freivalds verification.

\paragraph{Input-binding constraints.}
For each opened polynomial value $v$ used in the circuit (e.g., an entry 
of $u$ or $z$), the circuit enforces:
\[
  v_\text{circuit} = v_\text{open}.
\]
This prevents the prover from presenting inconsistent values between the
commitment verifier and the SNARK.

Such constraints can be implemented using the modular input-binding 
gadget of LegoSNARK~\cite{campanelli2019legosnark}.

\subsubsection{Polynomial-Opening Proofs}

The proof contains:
\[
  \pi_{\mathrm{PC}} = 
    \left\{
      \mathsf{Open}(f_{U_j}, i, U_{i,j}),
      \;
      \mathsf{Open}(f_{V_j}, i, V_{i,j})
    \right\}
\]
for all queried entries $U_{i,j}$ and $V_{i,j}$.

Each opening includes:
\begin{itemize}
  \item the evaluation value,
  \item a polynomial-opening proof (KZG or IPA),
  \item the index $i$.
\end{itemize}

\subsubsection{Full Proof Transcript}

A complete proof is the tuple:
\[
  \Pi =
  \left(
    y,\;
    \pi_{\mathrm{PC}},\;
    \pi_{\mathrm{SNARK}}
  \right).
\]

\paragraph{Size of the Proof.}
The proof size is:
\[
  O(\#\text{openings}) \;+\; O(1)
\]
where the SNARK proof is constant-sized (or $O(k)$ in some schemes).

Because the number of openings is proportional only to the sparsity or
structure of $(r, x)$, the total size is independent of $n$.
\subsection{Complexity Analysis}

We analyze the prover cost, verifier cost, proof size, and preprocessing
requirements of our protocol.  
All complexities are expressed in terms of the sequence length $n$ and
the low-rank parameter $k$ (where $k \ll n$).

\subsubsection{Prover Time}

\paragraph{Online computation.}
The prover computes
\[
  z = V^\top x \quad \text{and} \quad y = U z,
\]
which requires $O(nk)$ field operations.

\paragraph{Polynomial openings.}
To support Freivalds’ check, the prover opens $O(k)$ entries drawn from
the columns of $U$ and $V$.
For KZG commitments, each opening requires $O(1)$ prover time and
constant-size proofs.
For IPA-based commitments (e.g., DARK, Kim et al. 2023), each opening
requires $O(\log n)$ prover time.

Thus, the total cost of generating all openings is:
\[
  O(k) \text{ (KZG)} 
  \qquad\text{or}\qquad 
  O(k \log n) \text{ (IPA)}.
\]

\paragraph{SNARK proof generation.}
The SNARK circuit contains only $O(k)$ multiplications for computing
$u^\top z$.
Witness generation time is therefore:
\[
  O(k)
\]
plus the overhead of the underlying SNARK backend.

\paragraph{Total prover time.}
Overall, the prover cost is:
\[
  O(nk) + O(k \log n) + O(k) 
  = O(nk + k \log n).
\]
For KZG commitments, the term reduces to:
\[
  O(nk).
\]

\subsubsection{Verifier Time}

The verifier performs three tasks:

\paragraph{Sampling the challenge.}
Sampling a random $r \in \mathbb{F}^n$ requires negligible time;  
in practice the verifier can use sparse or structured randomness.

\paragraph{Checking polynomial openings.}
Each opening is verified in:
\[
  O(1) \text{ (KZG)} 
  \qquad\text{or}\qquad 
  O(\log n) \text{ (IPA)}.
\]

With $O(k)$ openings, the total verification cost is:
\[
  O(k) \text{ (KZG)} 
  \qquad\text{or}\qquad 
  O(k \log n) \text{ (IPA)}.
\]

\paragraph{Checking the SNARK proof.}
The SNARK proof verifies in:
\[
  O(1) \ \text{or}\ O(k),
\]
depending on the scheme.

\paragraph{Total verifier time.}
Thus, the verifier runs in:
\[
  O(k) \text{ (KZG)} 
  \qquad\text{or}\qquad 
  O(k \log n) \text{ (IPA)}.
\]

Importantly, the verifier’s cost is independent of the sequence length
$n$.

\subsubsection{Proof Size}

The full proof consists of $O(k)$ openings plus a constant-size SNARK
proof:
\[
  |\Pi| 
  = 
  O(k) \text{ (KZG)} 
  \qquad\text{or}\qquad 
  O(k \log n) \text{ (IPA)}.
\]

No component scales with $n$.

\subsubsection{Preprocessing Cost}

The setup phase commits to the $k$ columns of each matrix $U$ and $V$.
Each column requires a single polynomial commitment.
Thus preprocessing time and storage are:
\[
  O(nk) \text{ to encode polynomials,}
  \qquad 
  O(k) \text{ commitments (constant-size or logarithmic).}
\]

Setup happens once and is amortized across all proofs.

\subsubsection{Comparison With Naïve Approaches}

A naïve SNARK encoding of the attention operation $y = A x$ with
$A = U V^\top$ would require:

\[
  \Theta(n^2) \text{ multiplications in the circuit},
\]
since the dense $n \times n$ matrix $A$ would need to be materialized.

In contrast, our protocol reduces the verification to:
\[
  O(k) \text{ circuit operations} \qquad (k \ll n),
\]
and eliminates all dependence on $n$ in verifier time and proof size.

\paragraph{Summary.}
The final complexities of our protocol are:
\[
  \textbf{Prover: } O(nk + k \log n)
\]
\[
  \textbf{Verifier: } O(k) \text{ to } O(k \log n)
\]
\[
  \textbf{Proof Size: } O(k) \text{ to } O(k \log n)
\]
\[
  \textbf{Verifier independence from } n: \text{Yes}
\]
demonstrating that verification remains efficient even for extremely
large sequence lengths.

\section{Security Analysis}

\subsection{Completeness}

We show that if the prover follows the protocol honestly and the matrices
$U$ and $V$ committed in the preprocessing phase are correct, then the
verifier always accepts.

\paragraph{Correct polynomial openings.}
Since each column $U_j$ and $V_j$ is committed as a polynomial
$f_{U_j}, f_{V_j}$ during preprocessing, all openings produced by an
honest prover correspond exactly to evaluations of these polynomials.
Thus every polynomial-opening proof in $\pi_{\mathrm{PC}}$ verifies
successfully.

\paragraph{Correct low-rank computation.}
The prover computes
\[
  z = V^\top x,\qquad y = U z,
\]
using the same matrices $U$ and $V$ that were committed.
Therefore,
\[
  y = U(V^\top x)
\]
holds identically.

\paragraph{Correct Freivalds relation.}
For the verifier-chosen random vector $r$, the prover computes:
\[
  s_1 = r^\top y,
  \qquad
  u = r^\top U,
  \qquad
  s_2 = u^\top z.
\]
Since $y = U z$, it follows that:
\[
  s_1 = r^\top (U z)
      = (r^\top U) z
      = u^\top z
      = s_2.
\]
Thus the scalar equality enforced inside the SNARK holds.

\paragraph{Correct SNARK proof.}
Because all circuit constraints are satisfied (including the input-binding 
constraints tying the circuit inputs to the polynomial openings), the SNARK
prover produces a valid proof $\pi_{\mathrm{SNARK}}$.
Therefore, the SNARK verifier accepts.

\paragraph{Final completeness.}
The verifier accepts if and only if:
(i) all polynomial openings verify, and
(ii) the SNARK proof verifies.
Both conditions hold for an honest prover.

\[
  \textbf{Therefore, the protocol is complete.}
\]


\subsection{Soundness}

We show that no malicious prover can convince the verifier of an
incorrect output $y \neq U(V^\top x)$ except with negligible
probability.  
Soundness follows from four independent components:
(1) soundness of polynomial commitments,  
(2) soundness of SNARK input binding,  
(3) soundness of Freivalds’ randomized check, and  
(4) a final composition theorem.

\subsubsection{Soundness of Polynomial Openings}

Let $\mathsf{Commit}$ be the polynomial-commitment scheme (KZG or IPA).
By the binding property of the commitment scheme, a prover cannot open a
polynomial $f$ at an index $i$ to two different values.  
Thus, for every revealed entry $U_{i,j}$ or $V_{i,j}$, the value
provided by the prover must equal the evaluation of the polynomial
committed during setup.

Therefore, the prover cannot cheat by substituting incorrect entries in
$U$ or $V$ without producing an invalid opening proof.

\subsubsection{Soundness of SNARK Input Binding}

The SNARK includes input-binding constraints which enforce that every
scalar used inside the circuit (e.g., components of $u = r^\top U$ or
$z = V^\top x$) matches the corresponding value that appears in a
verified polynomial opening.

Thus, even if the prover could produce fake values for $u$ or $z$, the
SNARK would reject unless these values exactly match the ones associated
with the polynomial openings.  
This prevents *equivocation* between the commitment layer and the circuit.

\subsubsection{Soundness of the Freivalds Relation}

Suppose the prover outputs an incorrect vector $y' \neq U(V^\top x)$.
Let
\[
  s_1' = r^\top y', \qquad
  s_2 = (r^\top U)(V^\top x).
\]

Since $y' \neq U(V^\top x)$, the difference
\[
  \Delta = y' - U(V^\top x)
\]
is nonzero.
For a random $r \leftarrow \mathbb{F}^n$, Freivalds’ analysis shows:
\[
  \Pr_r[\,r^\top y' = (r^\top U)(V^\top x)\,]
  =
  \Pr_r[\,r^\top \Delta = 0\,]
  \le \frac{1}{|\mathbb{F}|}.
\]

Thus a cheating prover can satisfy the scalar equation only with
probability at most $1/|\mathbb{F}|$ per attempt, and this probability
drops exponentially with repeated checks.

\subsubsection{Final Soundness Theorem}

We combine the three components:

\begin{itemize}
  \item If the prover presents incorrect openings for $U$ or $V$, the
        polynomial-commitment verifier rejects.
  \item If the prover presents openings but supplies inconsistent values
        inside the SNARK, the input-binding constraints cause the SNARK
        to reject.
  \item If the prover supplies consistent openings but computes an
        incorrect output $y'$, Freivalds’ check rejects except with
        probability at most $1/|\mathbb{F}|$.
\end{itemize}

Therefore, for any malicious (polynomial-time) prover,
\[
  \Pr[\text{Verifier accepts incorrect } y']
  \le
  \frac{1}{|\mathbb{F}|} + \mathsf{negl}(\lambda).
\]

\[
  \boxed{
    \text{The protocol is sound under the security of the polynomial
    commitment scheme and the underlying SNARK.}
  }
\]
\section{Extensions to Multi-Head Attention}

Transformer architectures use multi-head attention, where $h$ parallel
attention heads operate on the same input.  
We show that our protocol extends naturally to the multi-head setting
without increasing verification cost beyond a factor linear in the
number of heads, and in many cases with only constant overhead.

\subsection{Multi-Head Structure}

A multi-head attention layer computes, for heads $1 \le t \le h$:
\[
  S^{(t)} = \frac{Q W_Q^{(t)} \, (K W_K^{(t)})^\top}{\sqrt{d}},
  \qquad
  A^{(t)} = \mathrm{softmax}(S^{(t)}),
\]
\[
  Y^{(t)} = A^{(t)} \, (V W_V^{(t)}),
\]
and then concatenates the head outputs:
\[
  Y = \mathrm{Concat}\big( Y^{(1)},\dots,Y^{(h)} \big) W_O.
\]

Following efficient multi-head models (e.g., Linformer, Performer), each
attention matrix $A^{(t)}$ is approximated using its own low-rank
factorization:
\[
  A^{(t)} \approx U^{(t)} (V^{(t)})^\top,
\]
with rank $k$ per head.

\subsection{Extending Our Commit-and-Prove Framework}

Our protocol extends to multi-head attention in a modular way:

\begin{itemize}
  \item For each head $t$, the model owner commits to the matrices
        $U^{(t)}$ and $V^{(t)}$ during preprocessing.
  \item For each $t$, the prover computes:
        \[
          z^{(t)} = (V^{(t)})^\top x, \qquad
          y^{(t)} = U^{(t)} z^{(t)}.
        \]
  \item For each $t$, the prover produces polynomial openings for the
        entries used in $r^\top U^{(t)}$ and $(V^{(t)})^\top x$.
  \item A separate Freivalds check is applied per head:
        \[
          r^{(t)\top} y^{(t)} = 
          \big(r^{(t)\top} U^{(t)}\big) \, (z^{(t)}).
        \]
  \item The SNARK includes $h$ parallel scalar-equality checks, one for
        each head.
\end{itemize}

\subsection{Circuit Size and Verifier Cost}

The SNARK circuit for a single head requires $O(k)$ steps.
For $h$ heads, the total circuit size is:
\[
  O(hk).
\]

Since $k \ll n$ and $h$ is typically small (e.g., $8$--$16$), the cost
remains extremely low compared to any naive encoding of self-attention.

Polynomial openings per head are independent, so the total number of
openings is:
\[
  O(hk).
\]

Verifier cost scales linearly with $h$:
\[
  O(hk) \quad \text{(KZG)} 
  \qquad\text{or}\qquad 
  O(hk \log n) \quad \text{(IPA)}.
\]

\subsection{Shared-Randomness Optimization}

To reduce overhead further, the verifier may reuse a shared challenge
vector $r$ across all heads:
\[
  r^{(1)} = \dots = r^{(h)} = r,
\]
which preserves soundness, as each head uses distinct committed matrices
$U^{(t)}, V^{(t)}$.

This reduces the cost of challenge sampling and may enable batching of
polynomial-opening proofs.

\subsection{Final Output Aggregation}

After verifying each head output $y^{(t)}$, the prover demonstrates that
the concatenated vector
\[
  Y = \mathrm{Concat}\big( y^{(1)}, \dots, y^{(h)} \big) W_O
\]
was computed correctly.

Since $W_O$ is a publicly known linear transformation, its verification
requires only a small additional SNARK gadget with $O(hd)$ cost, and
does not affect the asymptotic complexity.

\paragraph{Summary.}
Our protocol extends to multi-head attention with overhead linear in the
number of heads $h$ and independent of the sequence length $n$.
No new cryptographic components are required, and all guarantees
(completeness, soundness, succinctness) remain intact.

\section{Handling Softmax and Nonlinearities}

The standard self-attention mechanism applies a softmax normalization
to the similarity matrix:
\[
  A = \mathrm{softmax}\!\left(\frac{QK^\top}{\sqrt{d}}\right),
\]
where the softmax is taken row-wise.  
Softmax is non-linear and expensive to express inside succinct proof
systems, as it requires evaluating exponentials, divisions, and
normalization constants.  
We discuss how our verifiable linear-attention framework naturally
avoids the need to handle exact softmax inside the SNARK.

\subsection{Linear-Attention Variants}

A central observation is that many efficient Transformer variants
replace the softmax kernel with a linear or kernelized approximation:
\[
  A \approx \phi(Q)\, \phi(K)^\top,
\]
for a suitable feature map $\phi$.
Examples include:

\begin{itemize}
  \item Performer: $\phi(x)$ approximates $\exp(x)$ using random features.
  \item Linear Transformers: kernelized dot-product approximations.
  \item FAVOR+, Random Fourier features.
  \item Softmax-free (e.g., CosFormer, AFT), which completely remove 
        the exponential.
\end{itemize}

These systems eliminate the need to evaluate exponentials explicitly and
yield attention matrices that admit a low-rank factorization of the
form:
\[
  A \approx U V^\top.
\]

Our protocol is compatible with all such models, as its complexity
depends only on the low-rank structure of $U$ and $V$, not on the
presence or absence of softmax.

\subsection{Low-Rank Softmax Approximations}

Even in architectures that retain softmax, it is often observed that the
attention matrix $A = \mathrm{softmax}(QK^\top)$ is numerically low rank
for practical inputs.  
Methods such as Linformer, Nyströmformer, and Linearized LRA directly
approximate the softmax-attention matrix by constructing explicit
low-rank factors:
\[
  A \approx U V^\top
\]
without reproducing the exact softmax computation.

In such cases, the model designer provides $U$ and $V$ as part of the
model parameters or precomputation phase.  
Our protocol then verifies attention entirely through these linear
factors, without requiring the prover to evaluate softmax inside the
SNARK or polynomial-commitment layer.

\subsection{Why Exact Softmax Evaluation is Not Required}

There are two reasons we do not require exact softmax verification:

\paragraph{1. Soundness depends on the correctness of $U$ and $V$, not on softmax.}

Once the model owner commits to $U$ and $V$ during preprocessing, the
verifier only checks that:
\[
  y = U(V^\top x).
\]
Whether $U$ and $V$ approximate a softmax matrix or come from a
softmax-free architecture does not affect cryptographic correctness.

\paragraph{2. Exact softmax-based attention is infeasible in succinct proofs.}

Computing:
\[
  \exp(\cdot), \quad \sum_j \exp(\cdot_j), \quad \text{and the division}
\]
inside a SNARK would require thousands of constraints per entry and
potentially large rational approximations.  
Performing this for all $n$ rows of the attention matrix results in
$\Omega(n^2)$ constraints, defeating succinctness.

Thus, low-rank linear structures are the only viable path for scalable
verification.

\subsection{Verifying the Approximation Quality (Optional)}

If desired, the model owner can additionally commit to an error bound:
\[
  \| A - UV^\top \| \le \varepsilon,
\]
and publish $\varepsilon$ as an external guarantee.

Verifying such an error bound does not require a SNARK—they are model
quality guarantees analogous to those provided by Linformer and other
efficient Transformer literature.

Our protocol treats $U$ and $V$ as authoritative parameters and does not
need to prove their training correctness.

\subsection{Summary}

Our protocol avoids the need to implement softmax inside a SNARK by
building on linear-attention and low-rank approximations, which are
widely used in modern efficient Transformer architectures.  
This preserves both efficiency and cryptographic soundness while
remaining faithful to how high-performance self-attention models are
deployed in practice.

\section{Discussion and Limitations}

We discuss several important limitations of our approach and clarify the
scope under which our protocol provides meaningful guarantees.  
These limitations highlight fundamental challenges in verifiable
attention and indicate promising directions for future work.

\subsection{Dependence on Low-Rank Structure}

Our protocol assumes that the attention matrix admits (or is supplied
with) a low-rank factorization $A \approx UV^\top$ of rank $k \ll n$.
This is consistent with a large body of efficient-Transformer research,
but not all models satisfy this assumption.

\paragraph{High-rank attention.}
If $A$ is full-rank or numerically high-rank, then $k$ may be large,
leading to prover and verifier costs of:
\[
  O(nk) \quad \text{and} \quad O(k)
\]
which become impractical when $k \approx n$.

Thus, the protocol inherits the limitations of low-rank Transformer
approximations.

\paragraph{Quality of approximation.}
We do not prove that $UV^\top$ is a faithful approximation of the true
attention matrix.  
This is a model-level guarantee, not a cryptographic one.  
In particular, the verifier only checks:
\[
  y = U(V^\top x),
\]
not whether $y$ approximates the output of softmax attention.

\subsection{Preprocessing Assumptions}

Our scheme requires a one-time preprocessing phase where the model
owner commits to all columns of $U$ and $V$.

\paragraph{Trust in preprocessing.}
If preprocessing is malicious, the prover can always produce proofs
consistent with the committed (but incorrect) matrices.  
This is inherent to all preprocessing-based SNARK constructions and
commit-and-prove protocols.

Trusted setup or verifiable setup procedures (e.g., MPC-based generation
of commitments) can mitigate this risk but do not eliminate it entirely.

\subsection{Polynomial-Commitment Overheads}

While polynomial openings are asymptotically succinct, the constant
factors may be non-trivial in practice:

\begin{itemize}
  \item KZG requires pairings, which can dominate verifier time.
  \item IPA-based commitments (DARK or Kim et al. 2023) require 
        $O(\log n)$ group operations per opening.
  \item Large $k$ leads to $O(k)$ openings per proof.
\end{itemize}

These costs are unavoidable: opening a linear number of polynomial
entries is information-theoretically necessary for a single Freivalds
relation across $k$-dimensional intermediate representations.

\subsection{Circuit-Level Limitations}

Our SNARK circuit is extremely small—only $O(k)$ multiplications—but it
remains linear in $k$.  
This is unavoidable because:

\[
  s_2 = u^\top z
\]
is a $k$-dimensional inner product, and verifying it inside any
NP-succinct proof requires at least $\Omega(k)$ constraints.

We therefore inherit a linear dependence on $k$ in the size of the
SNARK.

\subsection{Randomness and Repetition}

Our soundness relies on Freivalds’ check, which fails with probability
$1/|\mathbb{F}|$.

\paragraph{Single-check limitations.}
While this probability is negligible for large fields (e.g., $|\mathbb{F}| = 2^{254}$),
it does require true uniform randomness.  
Deterministic or adversarially biased $r$ may invalidate soundness.

Multiple independent checks reduce failure probability exponentially but
also increase proof size and verifier time linearly.

\subsection{Nonlinear Components Outside the Scope}

Our protocol verifies only linear attention layers.
We do not verify:

\begin{itemize}
  \item activation functions (ReLU, GELU),
  \item softmax itself,
  \item normalization layers,
  \item MLP feed-forward layers.
\end{itemize}

These components require substantial additional circuit cost or
approximation and are orthogonal to the problem addressed by this work.

\subsection{Summary}

Our protocol provides efficient verification for low-rank and
softmax-free attention mechanisms but does not aim to verify full
Transformers or nonlinear components.  
Its efficiency depends critically on the low-rank assumption, the
soundness of polynomial-commitment schemes, and the correctness of the
preprocessing phase.

These limitations are aligned with the constraints of existing
verifiable ML systems and highlight the inherent complexity of scalable
verification of general attention mechanisms.

\section{Experimental Evaluation}
\label{sec:experiments}

In this section we evaluate the concrete performance of our protocol.
Our goal is not to compete with highly optimized zkML systems, but to
validate the asymptotic claims from Section~\ref{sec:construction} and
demonstrate that verification cost is independent of the sequence length
$n$ and scales linearly in the rank parameter $k$ and the number of
heads $h$.

\subsection{Implementation Setup}
\label{subsec:impl-setup}

We implemented a prototype of our protocol using a KZG polynomial
commitment scheme and a Groth16 SNARK over the BLS12-381 curve.
The implementation is written in \texttt{[language]} using
\texttt{[crypto library]}, and all experiments were run on a
\texttt{[CPU model]}, \texttt{[GHz]}, with \texttt{[RAM]} and
\texttt{[OS / version]}.

We generate synthetic low-rank matrices
$U, V \in \mathbb{F}^{n \times k}$ by sampling each entry uniformly at
random from the field and, optionally, applying a simple orthogonalization
step to avoid degenerate cases.
Input vectors $x \in \mathbb{F}^n$ are sampled uniformly at random.
Unless otherwise stated, results are averaged over $T$ independent
trials.

Table~\ref{tab:params} summarizes the default parameter choices used in
our experiments.

\begin{table}[t]
  \centering
  \begin{tabular}{l l}
    \toprule
    Parameter & Default value \\
    \midrule
    Field & BLS12-381 scalar field $\mathbb{F}_r$ \\
    Sequence length $n$ & $4{,}096$ (unless varied) \\
    Rank $k$ & $64$ (unless varied) \\
    Heads $h$ & $4$ (unless varied) \\
    Commitment scheme & KZG over BLS12-381 \\
    SNARK backend & Groth16 \\
    Trials per data point & $T = 20$ \\
    \bottomrule
  \end{tabular}
  \caption{Default experimental parameters.}
  \label{tab:params}
\end{table}

\subsection{Microbenchmarks: Prover and Verifier Cost}
\label{subsec:microbenchmarks}

We first study how the prover and verifier costs scale with the rank
parameter $k$.
We fix the sequence length to $n = [\text{e.g., }4{,}096]$ and vary
$k \in \{8, 16, 32, 64, 128\}$.

On the prover side, we decompose the total time into three components:
\begin{enumerate}
  \item $T_{\mathrm{matvec}}$: time to compute $z = V^\top x$ and
        $y = U z$,
  \item $T_{\mathrm{open}}$: time to generate KZG polynomial openings
        for the required entries of $U$ and $V$,
  \item $T_{\mathrm{snark}}$: time to construct the Groth16 proof for
        the Freivalds scalar relation.
\end{enumerate}

Figure~\ref{fig:prover-vs-k} reports these times as a function of $k$.

\begin{figure}[t]
  \centering
  % Replace with your actual plot file, e.g.:
  % \includegraphics[width=0.8\linewidth]{fig/prover-vs-k.pdf}
  \fbox{\parbox[c][4cm][c]{0.8\linewidth}{\centering
      Placeholder for prover time vs.\ $k$ plot}}
  \caption{Prover time vs.\ rank $k$ for fixed sequence length $n$.
  The total prover time is decomposed into the low-rank matvec
  computation, polynomial openings, and SNARK proof generation.}
  \label{fig:prover-vs-k}
\end{figure}

Similarly, we measure verifier time, decomposed into:
\begin{enumerate}
  \item $T_{\mathrm{open\text{-}verify}}$:
        time to verify all KZG openings,
  \item $T_{\mathrm{snark\text{-}verify}}$:
        time to verify the Groth16 proof.
\end{enumerate}

Figure~\ref{fig:verifier-vs-k} shows that verifier cost scales linearly
with $k$ and remains independent of $n$, confirming the analysis in
Section~\ref{subsec:complexity}.

\begin{figure}[t]
  \centering
  % \includegraphics[width=0.8\linewidth]{fig/verifier-vs-k.pdf}
  \fbox{\parbox[c][4cm][c]{0.8\linewidth}{\centering
      Placeholder for verifier time vs.\ $k$ plot}}
  \caption{Verifier time vs.\ rank $k$ for fixed sequence length $n$.
  The cost is dominated by KZG-opening verification and scales
  linearly with $k$.}
  \label{fig:verifier-vs-k}
\end{figure}

Table~\ref{tab:micro} summarizes representative measurements for a
single parameter setting (e.g., $n = 4{,}096$, $k = 64$).

\begin{table}[t]
  \centering
  \begin{tabular}{l r}
    \toprule
    Quantity & Time (ms) \\
    \midrule
    $T_{\mathrm{matvec}}$ & [XX.X] \\
    $T_{\mathrm{open}}$ & [XX.X] \\
    $T_{\mathrm{snark}}$ & [XX.X] \\
    \midrule
    $T_{\mathrm{prover}}$ (total) & [XX.X] \\
    \midrule
    $T_{\mathrm{open\text{-}verify}}$ & [X.X] \\
    $T_{\mathrm{snark\text{-}verify}}$ & [X.X] \\
    \midrule
    $T_{\mathrm{verifier}}$ (total) & [X.X] \\
    \bottomrule
  \end{tabular}
  \caption{Representative prover and verifier times for
  $n = 4{,}096$ and $k = 64$.}
  \label{tab:micro}
\end{table}

\subsection{Scaling with Sequence Length $n$}
\label{subsec:scaling-n}

Next, we fix the rank parameter (e.g., $k = 64$) and vary the sequence
length
\[
  n \in \{1{,}024, 2{,}048, 4{,}096, 8{,}192, 16{,}384\}.
\]
For each $n$, we measure total prover time, total verifier time, and
proof size.

Figure~\ref{fig:scaling-n} shows that prover time grows linearly with
$n$, as predicted by the $O(nk)$ matvec cost.
In contrast, verifier time and proof size remain essentially flat as $n$
increases, confirming that verification is independent of the sequence
length.

\begin{figure}[t]
  \centering
  % \includegraphics[width=0.8\linewidth]{fig/scaling-n.pdf}
  \fbox{\parbox[c][4cm][c]{0.8\linewidth}{\centering
      Placeholder for scaling vs.\ $n$ plot}}
  \caption{Prover and verifier time vs.\ sequence length $n$
  (with fixed rank $k$).
  Prover time grows linearly in $n$, while verifier time is independent
  of $n$.}
  \label{fig:scaling-n}
\end{figure}

\subsection{Multi-Head Attention}
\label{subsec:multihead-experiments}

We also evaluate the cost of verifying multi-head attention.
We fix $n$ and $k$ and vary the number of heads
\[
  h \in \{1, 2, 4, 8\}.
\]

For each $h$, we instantiate $h$ independent low-rank attention heads,
each with its own committed matrices $U^{(t)}, V^{(t)}$.
We measure total prover time, total verifier time, and total proof size
for verifying all heads in a batch.

Figure~\ref{fig:multihead} reports the results.
We observe that both prover and verifier cost scale linearly in $h$,
while remaining independent of $n$, in line with the analysis in
Section~\ref{sec:multihead}.

\begin{figure}[t]
  \centering
  % \includegraphics[width=0.8\linewidth]{fig/multihead.pdf}
  \fbox{\parbox[c][4cm][c]{0.8\linewidth}{\centering
      Placeholder for multi-head scaling plot}}
  \caption{Prover and verifier time vs.\ number of heads $h$.
  Both scale linearly in $h$, while remaining independent of
  the sequence length $n$.}
  \label{fig:multihead}
\end{figure}

\subsection{Proof Size}
\label{subsec:proof-size}

Finally, we report total proof size as a function of $k$ and $h$.
For each $(k, h)$, we record the size of:
(1) all KZG openings, and (2) the Groth16 proof.

Figure~\ref{fig:proof-size} shows that proof size grows linearly in
$k h$, reflecting the $O(kh)$ number of openings, while the SNARK proof
remains constant-sized.

\begin{figure}[t]
  \centering
  % \includegraphics[width=0.8\linewidth]{fig/proof-size.pdf}
  \fbox{\parbox[c][4cm][c]{0.8\linewidth}{\centering
      Placeholder for proof size vs.\ $k,h$ plot}}
  \caption{Total proof size as a function of rank $k$ and number of
  heads $h$. The size scales linearly with $kh$, while the SNARK proof
  component remains constant-sized.}
  \label{fig:proof-size}
\end{figure}

\subsection{Summary of Experimental Findings}
\label{subsec:exp-summary}

Across all experiments, our prototype behaves consistently with the
theoretical analysis:

\begin{itemize}
  \item Prover time scales as $O(nk)$, dominated by the low-rank matvec
        operations.
  \item Verifier time and proof size scale as $O(k)$ (and linearly in
        the number of heads $h$), and are independent of the sequence
        length $n$.
  \item Multi-head attention is handled by parallel composition with
        overhead linear in $h$, without affecting asymptotic succinctness.
\end{itemize}

These results confirm that our protocol is practically viable for
verifiable low-rank attention in realistic Transformer-style deployments.



\section{Related Work}

\subsection{Low-Rank and Efficient Transformer Architectures}

A large body of work has shown that the attention matrix in
Transformers often exhibits significant redundancy and can be
approximated by low-rank or kernelized structures.
Linformer~\cite{wang2020linformer} reduces the attention complexity from
$O(n^2)$ to $O(n)$ by projecting keys and values into a low-dimensional
subspace.  
Nyströmformer~\cite{xiong2021nystromformer} applies the Nyström method
to approximate softmax attention using a small set of landmark points.

Other architectures eliminate softmax entirely.  
Performer~\cite{choromanski2021performer} uses positive random features
to express softmax attention as a kernel product
$\phi(Q)\phi(K)^\top$, enabling linear-time computation.  
CosFormer and AFT~\cite{shi2021cosformer, zhu2021aft} replace the
softmax operation with alternative feature maps.

Our protocol builds on this line of work by assuming that each attention
head admits a factorization $A \approx UV^\top$ of moderate rank.  
Importantly, unlike these prior works, we do not contribute a new
approximation technique; rather, we make such approximations *verifiable*
through succinct proofs.

\subsection{Verifiable Machine Learning and ZK-ML}

Verifiable inference has been studied in several contexts.
Early work on general-purpose zkML~\cite{liu2021zkcnn}
encodes complete neural networks as arithmetic circuits, which leads to
very high prover overhead due to the quadratic cost of dense attention
layers.  
Recent systems~\cite{ezkl, lu2023tensorplonk} incorporate more
efficient proof backends and support CNNs and MLPs, but verifying
Transformers remains challenging due to the $O(n^2)$ structure of
attention.

A number of works propose high-level ML verifiability using hash-based
integrity arguments~\cite{cormode2011verifying}, but these typically do
not provide zero-knowledge or succinctness.

To our knowledge, no prior work has provided a succinct, 
low-rank–aware verification protocol for attention mechanisms.
Our approach is orthogonal to general zkML systems:  
we do not encode softmax or full attention into a circuit, but instead
validate the correctness of the *linearized* attention using
polynomial-commitment openings combined with a SNARK-based scalar
equation.

\subsection{Polynomial Commitment Schemes and SNARKs}

Polynomial commitments are essential for succinct verification.
KZG~\cite{kzg10} provides constant-sized proofs and $O(1)$ verification,
while inner-product-based schemes such as DARK~\cite{bunz2019dark} and the
work of Kim et al.~\cite{kim2023etpc} provide transparent
alternatives with $O(\log n)$ overhead.

SNARKs such as Groth16~\cite{groth16}, PLONK~\cite{plonk2019}, and
Marlin~\cite{marlin2020} provide succinct verification of arithmetic
circuits.  
LegoSNARK~\cite{campanelli2019legosnark} introduced modular input-binding, enabling
SNARKs to verify external committed values while preventing
equivocation.

Our construction combines polynomial openings with SNARK input-binding
to ensure that the prover cannot mix-and-match inconsistent matrix
entries.  
We leverage the strengths of commitments (binding $U,V$) and SNARKs
(ensuring consistency and enforcing Freivalds’ relation) while keeping
the circuit size $O(k)$.

\subsection{Randomized Checking and Freivalds Verification}

Freivalds’ algorithm~\cite{freivalds1979} provides a randomized
$O(n^2)$ check for verifying matrix multiplication, with soundness error
$1/|\mathbb{F}|$.  
It has been widely used in streaming, sublinear algorithms, and
verifiable computation.

Prior zkML systems typically do not use Freivalds, because attention
matrices must be explicitly represented inside the circuit, defeating
succinctness.

Our protocol applies Freivalds’ check in a new setting—verifying
low-rank attention—and shows that randomized checking combined with
commitment soundness yields succinct verification with complexity
independent of the sequence length $n$.

\subsection{Positioning of Our Work}

In summary:

\begin{itemize}
  \item Low-rank Transformers provide $UV^\top$ structure.
  \item Polynomial commitments bind $U$ and $V$.
  \item SNARKs enforce inner-product consistency.
  \item Freivalds supplies probabilistic verification.
\end{itemize}

To the best of our knowledge, no prior work combines these components to
produce a succinct and low-rank–aware verification protocol for
attention layers.
Our work bridges efficient Transformer approximations and modern
succinct proof systems to enable practical and scalable verification of
attention operations.

\section{Conclusion}

We presented the first succinct verification framework for low-rank
self-attention layers, a core component of modern Transformer models.
Our approach combines polynomial commitments, SNARK input-binding
mechanisms, and a Freivalds-style randomized check to verify
attention computations of the form $y = U(V^\top x)$ with complexity
independent of the sequence length~$n$.

The resulting protocol achieves prover time $O(nk)$, verifier time
$O(k)$ (or $O(k\log n)$ with transparent commitments), and proof size
$O(k)$.
These bounds align with the structure of efficient Transformer variants
such as Linformer, Performer, and Nyströmformer, and demonstrate that
verification can scale to long sequence lengths provided that the
attention matrix admits a moderate-rank factorization.

Our construction highlights a practical path toward verifiable
Transformer inference without encoding dense attention into arithmetic
circuits.
Future work includes extending the protocol to additional components of
Transformer architectures, improving multi-head batching, and exploring
verifiable approximations for softmax and other nonlinear operations.


\bibliographystyle{plain}
\bibliography{refs}

\end{document}
